% Part: propositional-logic
% Chapter: propositional-logic 
% Section: completeness

\documentclass[../../../include/open-logic-section]{subfiles}

\begin{document}

\olfileid{pl}{prp}{com}

\olsection{اكتمال منطق القضايا}

\begin{defn}
\ollabel{def:MaxCon}
تكون المجموعة~$\Gamma$ من~!!{formula}s \emph{متسقة على نحو أعظمي} إذا
كانت متسقة، وإذا كانت~$\Delta$ مجموعة متسقة تحقق~$\Gamma\subseteq\Delta$،
فإن~$\Gamma=\Delta$.
\end{defn}

\begin{lem}[لمّة الصدق]
  \ollabel{lem:truth}
لتكن~$\Gamma$ متسقة على نحو أعظمي؛ عندئذ:
\begin{enumerate}
\item \ollabel{lem:truth-1} $!A\in\Gamma$ إذا وفقط إذا كان
  $\lnot !A\notin\Gamma$؛
\item \ollabel{lem:truth-2} $\Gamma\Proves !A$ إذا وفقط إذا كان
  $!A\in\Gamma$؛
\item \ollabel{lem:truth-3} $!A\lif !B\in\Gamma$ إذا وفقط إذا كان
  إما~$!A\notin\Gamma$ أو~$!B\in\Gamma$.
\end{enumerate}
\end{lem}

\begin{proof}
\begin{enumerate}
\item افترض أن~$!A\in\Gamma$. يقتضي الاتساق أن
  $\lnot !A\notin\Gamma$. وبالعكس، افترض أن~$\lnot !A\notin\Gamma$
  وأن~$!A\notin\Gamma$. عندئذ لا تنتمي أي منهما إلى~$\Gamma$، وبحسب
  \olref[axd]{prop:phi} تكون إحدى المجموعتين~$\Gamma\cup\{!A\}$
  و~$\Gamma\cup\{\lnot !A\}$ متسقة وتمثل امتدادًا حقيقيًا لـ$\Gamma$،
  وهذا يناقض أعظميتها. ومن ثم~$!A\in\Gamma$.
\item إذا كان~$!A\in\Gamma$، فمن الواضح أن~$\Gamma\Proves !A$.
  ومن ناحية أخرى، افترض أن~$\Gamma\Proves !A$. فإذا كان
  $!A\notin\Gamma$، كان~$\lnot !A\in\Gamma$ بحسب
  \olref{lem:truth-1}. وعندئذ~$\Gamma\Proves !A$ و
  $\Gamma\Proves\lnot !A$، وهذا تناقض.
\item تمرين.
\end{enumerate}
\end{proof}

يبين الاتجاه من اليسار إلى اليمين في
\olref{lem:truth}\olref{lem:truth-2} أن~$\Gamma$ \emph{مغلقة استنباطيًا}؛
أي إنه إذا كان~$\Gamma\Proves !A$، فإن~$!A\in\Gamma$ لكل~$!A$.

\begin{prob}
  أكمل برهان~\olref{lem:truth}.
\end{prob}

\begin{thm}[الاكتمال]
\ollabel{thm:completeness}
إذا كانت~$\Gamma$ متسقة، فهي قابلة للإرضاء.
\end{thm}

\begin{proof}
لتكن~$!A_0,!A_1,\dots$ قائمة شاملة لجميع~!!{formula}s اللغة. نعرّف
تعاوديًا متتالية متزايدة من مجموعات~!!{formula}s
$\Gamma_0,\Gamma_1,\dots$ بوضع:
\begin{align*}
\Gamma_0 & = \Gamma\\
\Gamma_{n+1} &  =
\begin{cases}
  \Gamma_n \cup \{!A_n\} & \text{إذا كانت $\Gamma_n\cup\{!A_n\}$
    متسقة؛}\\
  \Gamma_n \cup \{\lnot !A_n\} & \text{فيما عدا ذلك.}
\end{cases}
\end{align*}
ثم نعرّف:
\[
\Gamma^*=\bigcup_{0\le n}\Gamma_n.
\]
ويمضي البرهان الآن بإثبات الحقائق الآتية بالترتيب:
\begin{enumerate}
\item المجموعة~$\Gamma_n$ متسقة لكل~$n$ (بالاستقراء على~$n$، مع استعمال
  \olref[axd]{prop:phi})؛
\item $\Gamma^*$ متسقة؛
\item $\Gamma^*$ أعظمية.
\end{enumerate}
ثم نعرّف تقييمًا~$v$ بوضع~$v(p_i)=\True$ إذا وفقط إذا كان
$p_i\in\Gamma^*$. ويُبيَّن بالاستقراء على~$!A$ أن الانتماء إلى
$\Gamma^*$ يطابق الصدق بحسب~$v$ أيضًا في~!!{sentence}s الأعقد:
$\overline{v}(!A)=\True$ إذا وفقط إذا كان~$!A\in\Gamma^*$. وعلى وجه
الخصوص، يحقق~$v$ المجموعة~$\Gamma^*$؛ ولأن~$\Gamma\subseteq\Gamma^*$،
تكون~$\Gamma$ قابلة للإرضاء أيضًا، كما هو مطلوب.
\end{proof}

\begin{cor}
إذا كان~$\Gamma\Entails !A$، فإن~$\Gamma\Proves !A$.
\end{cor}

\begin{proof}
إذا كان~$\Gamma\Proves/ !A$، كانت~$\Gamma\cup\{\lnot !A\}$ متسقة
بحسب~\olref[axd]{prop:prov-incons}. ومن ثم تكون
$\Gamma\cup\{\lnot !A\}$ قابلة للإرضاء بحسب مبرهنة الاكتمال، ويكون
$\Gamma\Entails/ !A$.
\end{proof}

\begin{prop}[مبرهنة التراص]
\ollabel{prop:compactness}
تكون~$\Gamma$ قابلة للإرضاء إذا وفقط إذا كانت كل مجموعة جزئية
\emph{منتهية}~$\Gamma_0$ منها قابلة للإرضاء.
\end{prop}

\begin{proof}
تكون~$\Gamma$ غير قابلة للإرضاء إذا وفقط إذا كانت غير متسقة، وهذا يكافئ
وجود مجموعة جزئية منتهية~$\Gamma_0$ من~$\Gamma$ غير متسقة، وهو يكافئ
وجود مجموعة جزئية منتهية~$\Gamma_0$ من~$\Gamma$ غير قابلة للإرضاء.
\end{proof}

\begin{cor}
$\Gamma\Entails !A$ إذا وفقط إذا وجدت مجموعة جزئية منتهية~$\Gamma_0$
من~$\Gamma$ تحقق~$\Gamma_0\Entails !A$.
\end{cor}

\end{document}
